\chapter{기어식 기계의 원리}\label{ch:50-gear-mechanics}

5부의 릴레이식 봄베에는 도는 것이 없었다. 로터의 회전을 접점으로 흉내 냈다(\ref{sec:relay-shifter}절). 이 부의
기어식 봄베는 전시의 최종 목표이고, 1940년대의 실물처럼 드럼이 실제로 돈다. 드럼이 돌면 풀어야 할 문제가
넷 생긴다.

\begin{enumerate}
\item 도는 드럼 안의 배선에 전기를 어떻게 넣고 빼는가(드럼 기판과 접점 보드).
\item 드럼 수십 개를 어떻게 함께, 정확한 각도로 돌리는가(기어열과 정렬).
\item 빠른 드럼이 한 바퀴 돌 때 가운데 드럼을 한 칸 넘기는 일, 곧 \emph{자리올림}을 기어만으로 어떻게 하는가.
\item 관람객이 드럼을 빼고 끼우는 동안 기계의 상태를 어떻게 지키는가(브레이크와 인덱스).
\end{enumerate}

이 장은 이 넷을 차례로 다룬다. 판정에 쓰는 시각(슬롯 원판)은 \ref{sec:t-gear}절에서 이미 보았으므로 짧게만
다시 짚는다. 치수는 저장소의 \code{hardware/gear/layout.md}, \code{concept.md}, 결정 0002에서 가져왔고, 문서가
``추정''이라고 적은 값은 여기서도 추정이라고 적는다. 시제품(G.1, G.2)에서 실측하기 전의 값이다.

\section{드럼: 손으로 끼우는 스크램블러}\label{sec:gm-drum}

\Why 에니그마의 로터는 고정된 배선이 든 원판이다. 원판의 한쪽 면에 26개의 접점이 있고 반대쪽 면에 26개가 있으며,
안의 전선이 한쪽의 $i$번 접점을 반대쪽의 $W[i]$번 접점에 잇는다. 원판이 $o$칸 돌아가 있으면 들어오는 글자는
$o$칸 밀린 자리에서 배선을 만나고, 나갈 때 다시 $o$칸 되돌려진다. \ref{sec:through}절의 \code{through} 함수가
``들어갈 때 더하고, 표를 읽고, 나올 때 뺀다''로 계산한 것이 바로 이 물리적 사실이다.

봄베의 드럼은 이 로터를 한 면으로 접은 것이다. 드럼의 앞면 하나에 동심원 네 줄로 26개씩, 모두 104개의 접점이
있다. 두 줄은 들어가는 길(입력 쪽 링과 반사판 쪽 링), 두 줄은 돌아오는 길이다. 같은 로터 배선이 드럼 안에 두
벌 들어 있는 셈이다~\cite{ellsbury1998bombe}. 접점을 한 면에 모은 까닭은 드럼을 기계 앞에서 손으로 끼우고 빼기
위해서다. 드럼 뒤쪽 한 면만 기계의 접점판에 닿으면 되므로 드럼은 축에 꽂기만 하면 된다.

\General 오래된 라디오의 \term{로터리 스위치}{rotary switch}가 같은 원리다. 손잡이를 돌리면 가운데의 와이퍼가
둘레에 늘어선 고정 접점 가운데 하나에 닿는다. 봄베의 드럼은 이것을 104개 동시에 하는 로터리 스위치이고,
접점 쌍 사이에 로터 배선이 끼어 있는 것이다.

\Ours 이 프로젝트의 드럼은 \emph{인쇄 회로 기판 한 장}이다(\code{hardware/gear/drum-pcb}). 지름 96~mm, 두께
1.6~mm, 4층이고 부품은 하나도 없다. 앞면(F.Cu)에 패드 104개가 반지름 22, 29, 36, 43~mm의 네 링에 있고, 안쪽 두
층이 패드끼리를 잇는다(그림~\ref{fig:gm-drum}). 링의 역할과 배선 규칙은 한 줄이다.

\begin{lstlisting}[numbers=none,caption={\code{drum-pcb/README.md}: 배선 규칙. $W$는 \code{data/wirings/enigma.txt}의 로터 배선표.},label={lst:gm-drumrule}]
    copy A:  pad 4<i>  <->  pad 3<W[i]>        (outer <-> 3rd ring)
    copy B:  pad 1<i>  <->  pad 2<W[i]>        (inner <-> 2nd ring)
\end{lstlisting}

링 4와 링 1이 입력 쪽, 링 3과 링 2가 반사판 쪽이다. 로터 I이면 $W[\mathrm{A}] = \mathrm{E}$이므로 패드 4A와 3E가,
1A와 2E가 이어진다. 돌아오는 전류는 링 2에서 링 1로 같은 배선을 거꾸로 지나므로 $W^{-1}$이 저절로 된다. 두 벌의
방향이 같아서 드럼은 위·가운데·아래 어느 자리에 꽂아도 같은 기판으로 된다.

\begin{figure}[ht]
\centering
\begin{tikzpicture}[x=1cm,y=1cm]
  % ---- (가) 패드 면
  \begin{scope}
    \def\s{0.055}
    \fill[cgreen!25] (0,0) circle (48*\s);
    \draw[cgreen!70!black] (0,0) circle (48*\s);
    \fill[white] (0,0) circle (4.3*\s);
    \draw[cgray,dashed] (0,0) circle (17*\s);
    \foreach \rr/\ww in {22/2.66,29/3.50,36/4.35,43/5.20} {
      \foreach \i in {0,...,25} {
        \pgfmathsetmacro{\a}{90-\i*360/26}
        \fill[corange!80,rotate=\a] ({\rr*\s-0.099},{-\ww*\s/2}) rectangle ++(0.198,{\ww*\s});
      }
    }
    \foreach \i/\L in {0/A,1/B,2/C,3/D,4/E,5/F,6/G,7/H,8/I,9/J,10/K,11/L,12/M,13/N,14/O,15/P,16/Q,17/R,18/S,19/T,20/U,21/V,22/W,23/X,24/Y,25/Z} {
      \pgfmathsetmacro{\a}{90-\i*360/26}
      \node[font=\tiny\ttfamily] at (\a:{52*\s}) {\L};
    }
    % 강조: 4A-3E (A벌), 1A-2E (B벌)
    \fill[cred] ({90}:{43*\s}) circle (2.2pt);
    \fill[cred] ({90-4*360/26}:{36*\s}) circle (2.2pt);
    \draw[cred,line width=1pt,dashed] (90:{43*\s}) .. controls (75:{46*\s}) and (45:{40*\s}) .. ({90-4*360/26}:{36*\s});
    \fill[cblue] (90:{22*\s}) circle (2.2pt);
    \fill[cblue] ({90-4*360/26}:{29*\s}) circle (2.2pt);
    \draw[cblue,line width=1pt,dashed] (90:{22*\s}) .. controls (70:{24*\s}) and (45:{26*\s}) .. ({90-4*360/26}:{29*\s});
    \node[font=\footnotesize\bfseries] at (0,-3.35) {(가) 패드 면 (접점 보드 쪽에서 봄)};
    \node[note,anchor=west,text=cred] at (3.1,1.7) {4A--3E\\(A벌, 가는 길)};
    \node[note,anchor=west,text=cblue] at (3.1,0.6) {1A--2E\\(B벌, 오는 길)};
    \node[note,anchor=west] at (3.1,-0.6) {점선 원 $r = 17$:\\인덱스 원판 자리};
  \end{scope}
  % ---- (나) 옆에서 본 접촉
  \begin{scope}[xshift=8.6cm,yshift=-0.4cm]
    \draw[shaft] (0,-1.8) -- (0,2.6);
    \fill[cgreen!25,draw=cgreen!70!black] (-2.6,1.2) rectangle (2.6,1.45);
    \foreach \x in {-2.36,-1.97,-1.6,-1.21,1.21,1.6,1.97,2.36} \fill[corange] (\x-0.08,1.15) rectangle (\x+0.08,1.2);
    \fill[cgreen!45,draw=cgreen!70!black] (-2.9,-0.6) rectangle (2.9,-0.4);
    \foreach \x in {-2.36,-1.97,-1.6,-1.21,1.21,1.6,1.97,2.36} {
      \fill[cgray!40] (\x-0.1,-0.4) rectangle (\x+0.1,0.75);
      \fill[corange!60!black] (\x,0.95) circle (0.11);
    }
    \draw[<->] (3.1,-0.4) -- node[right,font=\scriptsize] {9.0 mm} (3.1,1.06);
    \node[note,anchor=west] at (-2.6,1.85) {드럼 기판 (축과 함께 돈다)};
    \node[note,anchor=west] at (-2.9,-0.95) {접점 보드 (서 있다)};
    \node[note,anchor=north] at (0,-1.25) {구르는 공 끝의 스프링 접점 104개};
    \node[font=\footnotesize\bfseries] at (0,-2.95) {(나) 옆에서 본 접촉};
  \end{scope}
\end{tikzpicture}
\caption{드럼 기판(로터 I). (가) 패드 104개가 네 링에 있다. 글자는 패드 쪽에서 보아 시계 방향이고 A가 12시다.
빨간 점선과 파란 점선은 안쪽 층을 지나는 연결 두 개의 예다. (나) 기계 쪽의 스프링 접점이 아래에서 패드를 누른다.
접점은 서 있고 드럼이 그 위를 돈다.}
\label{fig:gm-drum}
\end{figure}

\paragraph{드럼이 $o$칸 돌아가 있으면.} 접점 보드의 접점은 26자리 $j = 0,\ldots,25$에 고정되어 있다. 드럼이 기준
표시에 글자 $o$를 보이고 있으면 자리 $j$의 접점은 패드 $(j + o) \bmod 26$에 닿는다. 그러면 입력 쪽 접점 $j$에서
들어간 전류는 패드 $j+o$로 들어가 반사판 쪽 패드 $W[j+o]$로 나오고, 그 패드는 접점 $W[j+o] - o$ 위에 있다. 곧
\[
  j \;\longmapsto\; W[j+o] - o \pmod{26}
\]
이고, Rust 코어의 로터 함수(\ref{sec:through}절)와 글자 하나 다르지 않다. 기판 생성기가 배선표에서 기판을 만들고,
따로 짠 검증기가 기판 파일에서 연결을 다시 읽어 배선표와 대조하는 것은 릴레이 보드와 같은 방식이다
(\ref{ch:boardverify}장).

\paragraph{왜 전선이 아니라 기판인가.} 1940년대의 드럼은 104개의 연결을 대부분 전선으로, 일부를 구리 ``8자'' 링크로
만들었다~\cite{bombe-rebuild-drums}. 드럼이 30--60개 필요한 이 프로젝트에서 손으로 납땜하면 오배선이 생기고
검사가 어렵다. 기판은 배선표에서 생성되므로 틀릴 곳이 생성기 한 곳뿐이다. 그 대가로 지켜야 할 규칙이 하나 생겼다.
패드 면의 가운데(반지름 18~mm 안)에는 비아도, 구리도, 인쇄도 두지 않는다. 그 자리에 축과 드럼을 맞추는 인덱스
원판(지름 34~mm, \ref{sec:gm-brake}절)이 얹히기 때문이다. 검증기 \code{verify\_drum.py}가 이 금지 영역을 검사한다.

\section{접점 보드: 서 있는 쪽에 스프링을}\label{sec:gm-contact}

\Why 도는 것과 서 있는 것 사이에 전기를 건네려면 한쪽은 단단한 면, 다른 쪽은 그 면을 누르는 스프링이어야 한다.
실물 봄베는 드럼에 브러시를, 기계 앞면에 평평한 접점판을 두었다~\cite{us6812report1944}. 이 프로젝트는 반대로 한다.
\emph{기계에 스프링 접점, 드럼에 패드}다. 관람객이 드럼을 만지고 떨어뜨릴 수 있으므로 드럼은 튼튼하고 단순한 평판이어야
하고, 닳는 부품(스프링 접점)은 기계 쪽에 두어 정비하는 사람이 바꾸게 한다. 1942년 미 해군의 봄베 계획도 휠에
세그먼트, 기계에 교체형 브러시였다~\cite{desch1942memo}.

\General 볼펜의 끝을 떠올리면 된다. 볼펜은 종이 위를 \emph{미끄러지지} 않고 작은 공이 \emph{구른다}. 미끄러지면
닳고 걸리지만 구르면 마찰이 작다. 흔한 스프링 핀(포고 핀)은 위아래로 눌리는 용도로만 수명이 적혀 있고, 옆으로
끌리는 용도에는 근거가 없다.

\Ours 접점은 Mill-Max의 Omniball 7845-0이다~\cite{millmax-7845}. 끝이 지름 3.96~mm의 구르는 공이고, 제조사가
``옆으로 미끄러지거나 구르는 접촉''용이라고 적은 계열이다(결정 0002의 3절). 자유 높이 9.75~mm에서 0.76~mm 눌러
9.0~mm에서 쓴다. 전체 행정은 1.52~mm여서 패드 면의 흔들림, 접점 길이의 공차, 보드와 기둥의 공차를 모두 더한
최악($\pm0.6$~mm)까지 받는다. 접점 하나가 약 72~gf로 누르므로 104개가 드럼을 약 73~N(공칭)으로 밀어낸다. 이 힘은
드럼을 축에 붙드는 노브가 받는다(\ref{sec:gm-brake}절).

접점 보드는 드럼 자리마다 한 장(110 × 110~mm)이고 세 자리가 같은 설계다. 처음 개념 설계는 열마다 세로로 긴 보드 한
장을 두려 했으나, 닳는 것은 빠른 행뿐이므로 자리마다 나누면 갈아 끼울 소모품이 3분의 1로 준다. 보드의 링마다 26핀
헤더가 하나씩 있고, 보드 사이는 26가닥 리본 케이블로 잇는다. 한 열(스크램블러 하나)은 리본 여덟 개로 이어진다
(그림~\ref{fig:gm-column}).

\begin{figure}[ht]
\centering
\begin{tikzpicture}[x=1cm,y=1cm,
    st/.style={draw,fill=clight,rounded corners=2pt,minimum width=34mm,minimum height=12mm,align=center,font=\footnotesize},
    pt/.style={font=\scriptsize\ttfamily,fill=white,inner sep=1pt}]
  \foreach \y/\n/\lab in {0/b/{아래 = 느린 드럼\\(에니그마 오른쪽 로터)},2.4/m/{가운데 드럼},4.8/t/{위 = 빠른 드럼\\(에니그마 왼쪽 로터)}} {
    \node[st] (\n) at (0,\y) {\lab};
    \node[pt] (\n4) at (-1.3,\y-0.6) {R4};
    \node[pt] (\n3) at (-1.3,\y+0.6) {R3};
    \node[pt] (\n2) at (1.3,\y+0.6) {R2};
    \node[pt] (\n1) at (1.3,\y-0.6) {R1};
  }
  \node[blko,font=\footnotesize] (rf) at (0,6.6) {반사판 카드 (B 또는 C)};
  \node[font=\footnotesize] (in) at (-1.3,-1.6) {\code{IN} 잭};
  \node[font=\footnotesize] (out) at (1.3,-1.6) {\code{OUT} 잭};
  \draw[->,cred,line width=1pt] (in) -- (b4);
  \draw[->,cred,line width=1pt] (b3) -- (m4);
  \draw[->,cred,line width=1pt] (m3) -- (t4);
  \draw[->,cred,line width=1pt] (t3) -- (t3 |- rf.south);
  \draw[->,cblue,line width=1pt] (rf.south -| t2) -- (t2);
  \draw[->,cblue,line width=1pt] (t1) -- (m2);
  \draw[->,cblue,line width=1pt] (m1) -- (b2);
  \draw[->,cblue,line width=1pt] (b1) -- (out);
  \node[note,anchor=west,align=left] at (2.6,2.4) {빨강: 가는 길(드럼 안에서 R4 $\to$ R3, 배선 $W$)\\
    파랑: 오는 길(드럼 안에서 R2 $\to$ R1, 배선 $W^{-1}$)\\
    화살표 하나 = 26가닥 리본 하나, 한 열에 여덟 개\\[2mm]
    R$r$ = 반지름 22 / 29 / 36 / 43~mm의 링 1--4};
\end{tikzpicture}
\caption{열 하나(스크램블러 하나)의 배선. 케이블은 아래(느린) 드럼 쪽에, 반사판은 위(빠른) 드럼 쪽에 있다.
에니그마의 전류 순서(입력 $\to$ 오른쪽 $\to$ 가운데 $\to$ 왼쪽 $\to$ 반사판)를 따른 것이다.}
\label{fig:gm-column}
\end{figure}

\paragraph{같은 순서가 반대로 보인다.} 드럼 기판은 글자를 \emph{패드 쪽에서 보아} 시계 방향으로 놓는다. 접점
보드는 그 패드 면을 마주 보므로, 같은 물리적 순서가 보드의 앞면에서는 \emph{반시계} 방향으로 보인다. 그래서 접점
보드 생성기는 접점 자리를 반시계 방향으로 번호 매긴다. 두 문서가 서로 다른 방향을 적고 있는데 둘 다 옳은 것이다.
거울에 비친 시계가 거꾸로 도는 것과 같다. 이런 방향의 혼동은 말로만 따져서는 틀리기 쉽다. 그래서 열 검증기
\code{verify\_column.py}는 접점 보드, 리본, 반사판 카드의 기판 파일을 실제로 이어 붙여 에니그마 스크램블러와 비교하고,
보드를 거울상으로 뒤집은 것을 음성 시험으로 넣는다. 로터 순서 60가지 × 17,576위치, 모두 1,054,560가지가 Rust 코어와
같았다(\code{contact-board/README.md}).

\section{정렬: 패드 위에 확실히 얹히는 각도}\label{sec:gm-align}

\Why 접점은 패드 위에 \emph{완전히} 얹혀 있을 때만 믿을 수 있다. 패드 가장자리에 걸친 순간에는 접촉이 튀거나
이웃 패드와의 틈에 빠질 수 있다. 그래서 판정의 시험 전압(\code{STROBE})은 접점이 패드 위에 완전히 얹힌 구간의
\emph{가운데}에서만 건다. 남은 양쪽이 기계의 오차를 받는 여유다.

수를 넣어 보자. 안쪽 링(반지름 22~mm)의 글자 간격은 13.846°이고 패드는 그 50\%인 6.92° 폭이다. 접점 끝이 패드에
닿는 자리의 폭을 약 1~mm(안쪽 링에서 2.6°)로 잡으면, 접점 끝 전체가 패드 위에 있는 \term{완전 접촉}{full contact}
구간은 약 4.3°다. 그 가운데 50\%를 \code{STROBE}의 창으로 쓰면 창은 약 2.15° 폭이고, 빠른 드럼이 패드 가운데에서
$\pm 1.1°$까지 어긋나도 창이 완전 접촉 구간 안에 남는다(그림~\ref{fig:gm-window}). 빠른 드럼 60~rpm에서 4.3°는
약 12~ms, 창은 6~ms다. 그림~\ref{fig:t-slots}의 ``완전 접촉 $-3 \ldots +9$~ms''와 ``\code{STROBE} 접점 면
$0 \ldots 6$~ms''가 이것이다.

\begin{figure}[ht]
\centering
\begin{tikzpicture}[x=1.15cm,y=1cm]
  \draw[->] (-5.2,0) -- (5.4,0) node[right,font=\scriptsize] {각도 (°)};
  \foreach \t in {-4,-2,0,2,4} \draw (\t,0.06) -- (\t,-0.06) node[below,font=\tiny] {\t};
  \fill[corange!70] (-3.46,0.25) rectangle (3.46,0.6);
  \node[font=\scriptsize] at (0,0.42) {패드 (6.92°)};
  % ghost contacts at the limits
  \foreach \c in {-2.16,2.16} {
    \fill[cgray!50] (\c-1.3,0.65) rectangle (\c+1.3,0.85);
  }
  \node[note,anchor=south] at (-2.16,0.85) {접점 끝 2.6°\\(완전 접촉의 왼쪽 끝)};
  \node[note,anchor=south] at (2.16,0.85) {(오른쪽 끝)};
  \fill[cgreenbg] (-2.16,-0.75) rectangle (2.16,-0.5);
  \node[font=\scriptsize,text=cgreen] at (0,-0.62) {완전 접촉 약 4.3° (60 rpm에서 12~ms)};
  \fill[cbluebg] (-1.08,-1.25) rectangle (1.08,-1.0);
  \node[font=\scriptsize,text=cblue] at (0,-1.12) {\code{STROBE} 창 = 가운데 50\%};
  \draw[<->,cred] (1.08,-1.45) -- node[below,font=\scriptsize,text=cred] {허용 $+1.1°$} (2.16,-1.45);
  \draw[<->,cred] (-2.16,-1.45) -- node[below,font=\scriptsize,text=cred] {$-1.1°$} (-1.08,-1.45);
\end{tikzpicture}
\caption{안쪽 링의 패드 하나를 펴서 그린 것. 완전 접촉 구간은 패드 폭에서 접점 끝의 폭을 뺀 것이고, 창을 그
가운데 절반에 두면 양쪽에 1.1°씩 여유가 남는다. 구르는 공 끝은 닿는 자국이 더 작아 완전 접촉이 약 6.5°로
넓을 것으로 보지만(추정) 실측 전에는 이 여유를 셈에 넣지 않는다.}
\label{fig:gm-window}
\end{figure}

\General 오차를 쌓아서 셈하는 것을 \term{공차 누적}{tolerance stack-up}이라 한다. 종이 한 장의 두께가 0.1~mm $\pm$
0.01이면 100장을 쌓은 높이의 오차는 얼마일까. 모든 장이 동시에 가장 두꺼울 때를 가정하면 $\pm 1$~mm(\term{단순 합}{worst
case})다. 각 장의 오차가 서로 무관하게 흩어진다고 보면 대부분 서로 상쇄되어 $\pm 0.01 \times \sqrt{100} = \pm 0.1$~mm
근처(\term{제곱합의 제곱근}{root sum square})다. 단순 합은 안전하지만 비관적이고, 제곱합은 현실적이지만 운에 기댄다.
안전이 걸린 곳은 단순 합으로 맞추고 제곱합은 참고로 본다.

\Ours 빠른 드럼의 정렬 예산은 표~\ref{tab:gm-budget}이다. 단순 합 1.02°로 허용 1.1° 안에 들어가지만 여유가 0.1°뿐이다.
\code{layout.md}는 이곳을 ``이 설계에서 가장 빠듯한 곳''이라 적고 G.2 시제품의 측정 항목으로 남겼다. 다행히 백래시와
비틀림은 마찰의 방향이 늘 같아 대부분 \emph{일정한 치우침}이고, 일정한 치우침은 조립할 때 각도를 맞추면서 뺄 수 있다.

\begin{table}[ht]
\centering
\caption{빠른 드럼의 정렬 예산(\code{layout.md} 3.6절). ``계산''은 문서의 산술, ``추정''은 시제품에서 확인할 값.}
\label{tab:gm-budget}
\small
\begin{tabular}{lrl}
\toprule
원인 & 최악 & 근거 \\
\midrule
라인 축의 비틀림(12열) & 0.24° & 계산 \\
마이터 기어의 백래시 0.05--0.10~mm & 0.23° & 계산 \\
인덱스(핀과 홈)와 드럼 조립 지그 & 0.10° & 추정 \\
접점의 위치 오차 $\pm$0.1~mm (반지름 22) & 0.26° & 계산 \\
접점 보드의 위치 핀 & 0.04° & 계산 \\
패드 패턴 & 0.05° & 추정 \\
슬롯 원판과 센서(보정 뒤) & 0.10° & 추정 \\
\midrule
합: 단순 합 / 제곱합의 제곱근 & 1.02° / 0.45° & 허용 $\pm$1.1° \\
\bottomrule
\end{tabular}
\end{table}

가운데와 느린 드럼은 사정이 다르다. 이 드럼들은 \code{STROBE} 동안 서 있으므로 창이 완전 접촉 구간의 어디에 있든
상관없고, 허용이 구간 전체의 절반인 $\pm 2.1°$다. 정적인 항목에 자리올림 잠금의 유격 0.42°(\ref{sec:gm-carry}절)를
더해도 단순 합 1.34°다.

\paragraph{왜 기어를 줄줄이 잇지 않는가.} 한 행의 드럼 여섯 개(12유닛이면 열두 개)를 함께 돌리는 가장 쉬운 방법은
드럼 축마다 평기어를 달고 사이에 아이들 기어를 끼워 줄줄이 잇는 것이다. 모든 기어가 보여 전시물답다. 그러나 기어가
맞물리는 곳마다 백래시(이 사이의 틈)가 있고, 줄의 끝으로 갈수록 그것이 쌓인다. 맞물림 하나에 약 0.19°(백래시
0.1~mm, 피치 반지름 30~mm, 계산), 열 하나에 맞물림이 둘이므로 여섯 열이면 약 1.9°로 허용 1.1°를 넘는다. 그래서 행마다
기계 폭 방향으로 긴 지름 20~mm의 \term{라인 축}{line shaft}을 하나 두고, 드럼 축마다 1:1 \term{마이터 기어}{mitre
gear}(직각으로 맞물리는 베벨 기어)로 직접 받는다. 모든 드럼이 같은 축에서 받으므로 백래시가 쌓이지 않는다.
양동이를 사람에서 사람으로 건네면 흘리는 물이 쌓이지만, 모두가 같은 수도관에서 받으면 쌓이지 않는 것과 같다.

긴 축은 대신 비튼다. 축 끝에 토크 $T$가 걸리면 길이 $L$의 축은 $\theta = TL/(GJ)$만큼 비틀린다($G$는 강의 전단 탄성
계수, $J$는 단면의 극관성 모멘트로 지름의 네제곱에 비례한다). 드럼마다 마찰 토크가 고르게 걸리므로 끝 열의 비틀림은
그 절반인 $TL/(2GJ)$다. 12열(드럼당 0.6~N·m, 1.44~m)에서 지름 20~mm는 0.24°, 12~mm였다면 1.8°다(계산). 축을 굵게 한
이유다. 그리고 6유닛 단계에서는 구동부에 \emph{가까운} 7--12번 자리부터 채운다. 토크가 축의 짧은 쪽만 지나므로
비틀림이 0.06°로 준다(결정 0002의 6절). 왼쪽 절반이 비어 대칭은 깨지지만, 예산이 빠듯한 곳에서는 겉모습보다 예산을
택했다.

\section{기어열: 모터에서 세 줄의 축까지}\label{sec:gm-train}

\Why 봄베의 세 행은 서로 다른 빠르기로 돈다. 위의 빠른 행은 쉬지 않고 돌고(에니그마의 왼쪽 로터), 가운데 행은 빠른
행이 한 바퀴 돌 때 한 칸, 아래의 느린 행은 가운데 행이 한 바퀴 돌 때 한 칸 간다. 모터는 하나이므로 이 관계를
모두 기어가 만든다. 기어 두 개가 맞물리면 이의 수의 비가 회전수의 반비례가 된다. 이가 26개인 기어가 한 바퀴 돌면
이가 30개인 상대 기어는 26/30바퀴 돈다.

\General 자전거의 변속이 같은 원리다. 앞 체인링 이 40개, 뒤 스프로킷 이 20개이면 페달 한 바퀴에 바퀴가 두 바퀴
돈다. 두 기어 사이에 기어를 하나 더 끼우면(\term{아이들 기어}{idler gear}) 비는 그대로이고 도는 방향만 바뀐다.
아이들 기어의 이 수는 앞에서 곱해지고 뒤에서 나뉘어 지워지기 때문이다.

\Ours 구동부(\code{layout.md} 3.7절)에는 기계 폭 방향의 축이 다섯 개 있다. 빠른 라인 축 F, 아이들 축 I, 주기 축 C,
가운데 라인 축 M, 느린 라인 축 S다. 모터(폐루프 스테퍼)는 타이밍 벨트(HTD 5M, 18잇수 : 72잇수, 4:1)로 F를 돌린다.
그 뒤는 이렇다(그림~\ref{fig:gm-train}).

\begin{itemize}
\item F의 26잇수 $\to$ I의 26잇수 $\to$ C의 30잇수(모듈 1.5). F:C = 26:30이다. C가 한 바퀴 도는 동안 빠른 드럼은
30칸을 간다. 빠른 드럼 60~rpm에서 C는 52~rpm이고, 한 위치(빠른 드럼 한 칸)는 38.46~ms다.
\item C $\to$ M: 자리올림 1. C의 한 이빨 기어가 M의 26잇수 기어를 C 한 바퀴에 한 이씩 넘긴다(\ref{sec:gm-carry}절).
\item M $\to$ S: 자리올림 2. M의 한 이빨 기어가 S를 M 한 바퀴에 한 이씩 넘긴다.
\item 세 라인 축 $\to$ 드럼 축: 마이터 기어 25:25(1:1).
\end{itemize}

\begin{figure}[ht]
\centering
\begin{tikzpicture}[x=1cm,y=1cm]
  \def\s{0.021}
  \newcommand{\gmaxes}{\foreach \yy in {120,81,39,0,-120} {\fill[cgray] (0,\yy*\s) circle (1.6pt);}}
  % (가) 26:30
  \begin{scope}
    \gmaxes
    \foreach \yy/\rr in {120/19.5,81/19.5,39/22.5} \draw[cblue,fill=cbluebg,fill opacity=0.6] (0,\yy*\s) circle (\rr*\s);
    \foreach \nn/\yy/\tz in {F/120/26,I/81/26,C/39/30} \node[font=\scriptsize,anchor=west] at (0.55,\yy*\s) {\nn: \tz잇수};
    \node[font=\scriptsize,anchor=west,text=cgray] at (0.55,0) {M};
    \node[font=\scriptsize,anchor=west,text=cgray] at (0.55,-120*\s) {S};
    \node[font=\footnotesize\bfseries] at (0.6,-3.75) {(가) 26 : 26 : 30};
  \end{scope}
  % (나) carry 1
  \begin{scope}[xshift=4.6cm]
    \gmaxes
    \draw[corange,fill=corangebg,fill opacity=0.6] (0,0) circle (26*\s);
    \draw[corange,line width=0.8pt] (0,39*\s) circle (13*\s);
    \fill[corange] (-0.06,{26*\s}) rectangle (0.06,{26*\s-0.1});
    \node[font=\scriptsize,anchor=west] at (0.55,39*\s) {C: 한 이빨(13잇수 상당)};
    \node[font=\scriptsize,anchor=west] at (0.62,-0.15) {M: 26잇수};
    \node[font=\scriptsize,anchor=west,text=cgray] at (0.55,120*\s) {F};
    \node[font=\scriptsize,anchor=west,text=cgray] at (0.55,-120*\s) {S};
    \node[font=\footnotesize\bfseries] at (0.6,-3.75) {(나) 자리올림 1: C $\to$ M};
  \end{scope}
  % (다) carry 2
  \begin{scope}[xshift=10.0cm]
    \gmaxes
    \draw[cgreen,line width=0.8pt,dashed] (0,0) circle (81*\s);
    \fill[cgreen] (-0.07,{-81*\s}) rectangle (0.07,{-81*\s-0.1});
    \draw[cgreen,fill=cgreenbg,fill opacity=0.6] (0,-120*\s) circle (39*\s);
    \node[font=\scriptsize,anchor=west,align=left] at (0.1,0.45) {M: 한 이빨\\(54잇수 상당)};
    \node[font=\scriptsize,anchor=west] at (0.9,-120*\s) {S: 26잇수};
    \node[font=\scriptsize,anchor=west,text=cgray] at (0.55,120*\s) {F};
    \node[font=\footnotesize\bfseries] at (0.6,-3.75) {(다) 자리올림 2: M $\to$ S};
  \end{scope}
  \draw[<->] (-1.0,120*\s) -- node[left,font=\scriptsize] {120} (-1.0,0);
  \draw[<->] (-1.0,0) -- node[left,font=\scriptsize] {120} (-1.0,-120*\s);
\end{tikzpicture}
\caption{구동부의 축 다섯 개를 옆에서 본 것(축척, 원은 피치원). 다섯 축은 한 수직선 위에 있고 F--I--C--M의 축간
거리가 $39 + 42 + 39 = 120$~mm로 행 피치와 같다. 세 그림은 같은 축을 서로 다른 평면(기계 폭 방향의 다른 자리)에서
본 것이다. 한 이빨 기어의 이는 위치를 나타낸 작은 사각형이다.}
\label{fig:gm-train}
\end{figure}

축의 자리에 작은 산수가 있다. 기어 한 쌍의 축간 거리는 모듈 × (이 수의 합) ÷ 2다. F--I는 $1.5 \times 52 / 2 = 39$,
I--C는 $1.5 \times 56 / 2 = 42$, C--M은 자리올림 1의 기어(모듈 2, 13 + 26)로 39다. 합이 120~mm로 행 사이의 거리와
같아서 F와 M이 각각 빠른 행과 가운데 행의 라인 축과 같은 높이에 온다. 구동부의 다섯 축이 라인 축의 평면 하나에 들어가는
것이다. 아이들 I가 있는 이유도 여기 있다. F와 C를 직접 물리면 축간 거리가 $1.5 \times 56 / 2 = 42$로 고정되어 높이가
맞지 않는다. I가 끼면 거리를 39와 42로 나눌 수 있다. 방향은 F($+$) $\to$ I($-$) $\to$ C($+$) $\to$ M($-$) $\to$ S($+$)로
번갈아 바뀌고, 가운데 행만 라인 축의 마이터를 드럼 축의 반대쪽에 달아 세 행의 드럼이 모두 관람객 쪽에서 보아
시계 방향으로 가게 한다.

\paragraph{모터의 크기.} 드럼 하나의 마찰 토크는 0.2--0.6~N·m로 추정한다(접점 104개의 힘 × 마찰 계수 × 평균 반지름).
빠른 행 12개면 F에 2.4--7.2~N·m가 연속으로 걸린다. 가장 힘든 순간은 676위치에 한 번, 두 자리올림이 겹칠 때다. 느린
행을 넘기는 토크가 가운데 축을 거쳐 F로 오고, 자리올림 2는 느린 행을 가운데 행보다 54/26 = 2.08배 빠르게 돌리므로
\[
  7.2 + \bigl(7.2 + 7.2 \times \tfrac{54}{26}\bigr) \times \tfrac{13}{26} \times \tfrac{26}{30} \approx 16.8~\text{N·m}
\]
가 F에 걸린다(최악, 계산). 처음 개념 설계는 이 순간을 빼고 셌다. 결정 0002는 이 값으로 모터를 86~mm 각(NEMA 34)
8.5~N·m급 폐루프 스테퍼로 키우고, 4:1 벨트, 모터 전용 48~V 전원으로 정했다~\cite{leadshine-csm23485}. 같은 순간에
자리올림 1의 이에는 약 852~N이 걸리므로 그 기어는 황동이 아니라 강(S45C)으로 한다.

\section{자리올림: 한 이빨 기어와 잠금 원호}\label{sec:gm-carry}

\Why 자동차의 기계식 주행 거리계를 생각하자. 각 자리의 숫자 바퀴에는 이가 하나뿐인 작은 기어가 붙어 있다. 일의
자리 바퀴가 9에서 0으로 넘어가는 동안에만 그 이빨이 십의 자리 바퀴를 한 칸 밀고, 나머지 동안에는 십의 자리 바퀴를
건드리지 않는다. 회전 운동을 ``가끔 한 칸''이라는 운동으로 바꾸는 것이다. 봄베의 자리올림이 바로 이것이다.
빠른 드럼이 한 바퀴 돌 때마다 가운데 드럼이 한 칸, 가운데 드럼이 한 바퀴 돌 때마다 느린 드럼이 한 칸이다.

그런데 주행 거리계와 다른 점이 하나 있다. 빠른 드럼은 \emph{멈추지 않는다}. 가운데 드럼을 한 칸 넘기는 데는 시간이
들고, 그동안 가운데 드럼의 접점은 패드 사이에 걸려 있어 판정을 할 수 없다. 빠른 드럼의 이웃한 두 글자 사이의 짧은
틈에 자리올림과 판정을 함께 끼워 넣을 수는 없다. 실물 봄베의 해법은 한 주기를 26칸보다 길게 잡는 것이었다.
30칸 가운데 26칸에서 판정하고 4칸 동안 자리올림을 하는 ``30-point'' 기종, 26칸 + 13칸의 ``39-point'' 기종이
있었다~\cite{jackson-bh09,jackson-bh11,us6812report1944}. 이 프로젝트는 30칸을 택했다.

\General 영화 영사기는 필름을 한 칸씩 끊어 보내야 한다. 그 일을 하는 \term{제네바 기구}{Geneva drive}는 핀이 하나
달린 원판이 홈이 파인 십자 바퀴를 한 바퀴에 한 칸씩 돌리고, 나머지 동안에는 원판의 볼록한 원호가 십자 바퀴의
오목한 면에 얹혀 바퀴가 돌지 못하게 잠근다. 기계식 시계의 날짜판도 비슷하다. \emph{넘기는 부분}과 \emph{잠그는
부분}이 한 몸이라는 것이 핵심이다.

\Ours 주기 축 C의 둘레를 30칸으로 나누면 한 칸이 12°다. 칸 $k = 0, \ldots, 25$는 판정하는 칸이고(슬롯 원판에 슬롯이
있다), $k = 26, \ldots, 29$의 네 칸(48°)이 자리올림 구간이다(그림~\ref{fig:gm-cycle}). C에는 이빨이 하나뿐인 기어가
있다. 이빨은 13잇수 기어(피치원 지름 26~mm)의 이 하나에 해당하고, 둘레의 나머지는 반지름 11.95~mm의 \term{잠금
원호}{locking arc}다. M에는 26잇수 기어와 그 옆에 오목한 자리 26개가 파인 \term{잠금 원판}{lock disc}(반지름 27~mm)이
있다(그림~\ref{fig:gm-carry}).

\begin{itemize}
\item 이빨이 물리는 동안 C가 $360°/13 = 27.7°$ 돌고, 그동안 M이 이 하나, 곧 $360°/26 = 13.846°$를 간다. 27.7°는
자리올림 구간 48° 안에 넉넉히 들어간다. 맞물림을 빈 네 칸의 한가운데인 330°에 두면 316.2°--343.8°다.
\item 이빨이 물려 있지 않은 동안에는 C의 잠금 원호가 M의 오목한 자리에 얹혀 M이 돌지 못한다. 그래서 가운데 드럼은
판정하는 26칸 내내 정확한 자리에 서 있다.
\end{itemize}

\begin{figure}[ht]
\centering
\begin{tikzpicture}[x=1cm,y=1cm]
  \def\s{0.06}
  % driven: lock disc R27 with 26 pockets, centre (0,0)
  \fill[cgreenbg,draw=cgreen] (0,0) circle (27*\s);
  \foreach \i in {0,...,25} \fill[white] ({90+\i*360/26}:{27*\s}) circle (2.8*\s);
  \foreach \i in {0,...,25} \draw[cgreen] ({90+\i*360/26}:{27*\s}) ++({90+\i*360/26+90}:{2.8*\s}) arc ({90+\i*360/26+90}:{90+\i*360/26+270}:{2.8*\s});
  \fill[white,draw=cgray] (0,0) circle (10*\s);
  \node[font=\scriptsize] at (0,0) {M};
  % driver: arc disc R_a ~12 at (0, 39) with one tooth region
  \begin{scope}[shift={(0,{39*\s})}]
    \draw[corange,line width=1.2pt,fill=corangebg] (120:{12*\s}) arc (120:420:{12*\s}) -- (0,0) -- cycle;
    \fill[corange] (-0.12,{12*\s-0.05}) -- (0.12,{12*\s-0.05}) -- (0.08,{12*\s+0.37}) -- (-0.08,{12*\s+0.37}) -- cycle;
    \node[font=\scriptsize] at (0,-0.2) {C};
  \end{scope}
  \node[note,anchor=west,text=corange,align=left] at (1.2,2.75) {한 이빨: 맞물리는 동안\\M을 13.846° 넘긴다};
  \node[note,anchor=west,text=corange,align=left] at (1.2,1.95) {잠금 원호(반지름 11.95):\\나머지 동안 M을 붙든다};
  \node[note,anchor=west,text=cgreen,align=left] at (1.9,0.2) {잠금 원판(반지름 27)\\오목한 자리 26개};
  \node[note,anchor=west,align=left] at (1.9,-1.1) {틈 0.05~mm $\Rightarrow$ 유격 $\pm$0.33°\\(자리올림 2는 $\pm$0.42°)};
  % odometer sketch on the left
  \begin{scope}[xshift=-6.2cm,yshift=0.6cm]
    \foreach \x/\d in {0/3,1.1/9} {
      \draw[fill=white] (\x,0) rectangle ++(0.9,1.3);
      \node[font=\Large\ttfamily] at (\x+0.45,0.65) {\d};
    }
    \node[note] at (0.45,-0.35) {십의 자리}; \node[note] at (1.55,-0.35) {일의 자리};
    \draw[->,corange,line width=1pt] (1.35,1.55) to[bend right=40] node[above,font=\scriptsize,text=corange] {9$\to$0일 때만 한 칸} (0.55,1.55);
    \node[font=\footnotesize\bfseries] at (1.0,-1.2) {주행 거리계};
  \end{scope}
\end{tikzpicture}
\caption{자리올림의 원리(개념도, 축척은 대략). 왼쪽: 주행 거리계는 아랫자리가 넘어갈 때만 윗자리를 한 칸 민다.
오른쪽: 자리올림 1. 실제로는 이빨과 26잇수 기어가 한 평면에, 잠금 원호와 잠금 원판이 그 옆의 다른 평면에 있다.}
\label{fig:gm-carry}
\end{figure}

\paragraph{둘째 자리올림은 왜 큰가.} 같은 방법으로 M이 S를 넘기면 될 것 같지만 문제가 있다. M은 연속으로 돌지 않고
자리올림 1 때마다 13.846°씩만 걷는다. S를 넘기는 이빨은 M의 \emph{한 걸음 안에서} 일을 마쳐야 한다. 13잇수 상당의
이빨이면 M이 27.7° 돌아야 하므로 M의 두 걸음에 걸친다. 걸음 사이에 M이 서 있는 동안 S가 반쯤 넘어간 채 멈춰
버린다. 그래서 자리올림 2의 이빨은 54잇수 상당(모듈 3, 피치원 지름 162~mm)으로 크게 만든다. M이 $360°/54 = 6.67°$
도는 동안 S가 한 이를 가므로 한 걸음 안에 끝난다. 그 대가로 S는 M보다 2.08배 빠르게 돌려지고, 그것이
\ref{sec:gm-train}절의 16.8~N·m가 나온 이유다. 축간 거리는 $3 \times (54 + 26)/2 = 120$~mm로 역시 행 피치와 같다.

\paragraph{보조 디텐트를 달지 않는 이유.} 개념 설계는 넘겨진 행에 26칸 스프링 \term{디텐트}{detent}(볼펜을 딸깍
고정하듯 정해진 자리에 스프링이 걸리는 장치)를 달아 정위치를 잡으려 했다. 계산해 보니 드럼 12개의 마찰(2.4--7.2~N·m)을
이기고 중심을 잡으려면 디텐트가 반지름 30~mm에서 80--240~N으로 눌러야 했다. 그보다 약한 디텐트는 위치를 정하지
못한다. 거꾸로 마찰이 브레이크 노릇을 한다. 넘긴 뒤 관성으로 더 가는 각도는 운동 에너지 ÷ 마찰 토크로 가운데 행
0.09--0.26°, 느린 행 0.4--1.15°이고(계산), 그 뒤를 잠금 원호가 받는다. 그래서 결정 0002의 8절은 디텐트를 빼고 위치를
잠금 원호에 맡겼다. 원호와 원판 사이의 틈은 반지름 방향 0.05~mm이고, 이 틈이 허락하는 흔들림(유격)이 $\pm0.33°$
(자리올림 1), $\pm0.42°$(자리올림 2)다. 마찰의 방향이 늘 같으므로 실제로 서는 자리는 유격의 한쪽으로 일정하게
치우치고, 그 치우침은 조립 때 뺀다.

\paragraph{30칸 주기가 만드는 순서.} 주기 축이 한 바퀴 도는 동안 F는 30칸을 간다. 그래서 주기 축의 $Q$번째 회전,
$k$번째 칸에서 빠른 드럼이 보이는 글자는 $k$가 아니라
\[
  l = (30Q + k) \bmod 26 = (4m + k) \bmod 26 \qquad (m = Q \bmod 26)
\]
이다. 주기마다 빠른 드럼의 시작 글자가 네 칸씩 밀린다(그림~\ref{fig:gm-cycle} 아래). 한 주기 안에서는 $k$가 0부터 25까지
이어지므로 $l$도 26글자를 한 번씩 모두 지나고, 676번의 회전에서 17,576위치가 하나도 빠짐없이 한 번씩 시험된다. 다만
시험하는 \emph{순서}가 단순한 26진 계수가 아니다. 실물의 30-point 기종도 그랬다~\cite{jackson-bh09}. 이 사실은 제어
보드가 위치를 세고, 마지막 위치를 알고, 처음 자리로 돌아가는 방법을 모두 바꾼다(\ref{ch:52-gear-control}장). 기계가
시험하는 순서는 $n = 26Q + k$로 세고, 기계 밖으로 내보내는 위치 값은 드럼의 실제 글자 $(l, m, r)$이다.

\begin{figure}[ht]
\centering
\begin{tikzpicture}[x=1cm,y=1cm]
  \def\R{2.0}
  \foreach \k in {0,...,29} {
    \pgfmathsetmacro{\a}{90-12*\k}
    \pgfmathsetmacro{\b}{90-12*\k-6}
    \pgfmathsetmacro{\c}{90-12*\k+6}
    \ifnum\k<26
      \fill[cbluebg,draw=cblue!60] (0,0) -- (\b:\R) arc (\b:\c:\R) -- cycle;
    \else
      \fill[corangebg,draw=corange] (0,0) -- (\b:\R) arc (\b:\c:\R) -- cycle;
    \fi
  }
  \fill[white] (0,0) circle (0.9);
  \node[font=\scriptsize,align=center] at (0,0) {주기 축\\한 바퀴\\= 30칸};
  \foreach \k in {0,1,2,25,26,27,28,29} {
    \pgfmathsetmacro{\a}{90-12*\k}
    \node[font=\tiny] at (\a:{\R-0.35}) {\k};
  }
  \draw[corange,line width=1.6pt] ({90-316.2}:{\R+0.18}) arc ({90-316.2}:{90-343.8}:{\R+0.18});
  \node[font=\scriptsize,text=corange,anchor=south east] at ({90-330}:{\R+0.3}) {자리올림 1: 316°--344°};
  \node[font=\scriptsize,text=cblue,anchor=west] at (-15:{\R+0.25}) {$k = 0 \ldots 25$: 판정 (슬롯 있음)};
  % strip of fast drum letters
  \begin{scope}[xshift=-5.6cm,yshift=-3.4cm]
    \def\w{0.375}
    \node[font=\scriptsize,anchor=east] at (-0.1,0.45) {$k$};
    \node[font=\scriptsize,anchor=east] at (-0.1,0) {$Q = 0$};
    \node[font=\scriptsize,anchor=east] at (-0.1,-0.45) {$Q = 1$};
    \foreach \k in {0,...,29} {
      \pgfmathtruncatemacro{\ca}{mod(\k,26)+65}
      \pgfmathtruncatemacro{\cb}{mod(\k+4,26)+65}
      \ifnum\k<26 \def\fc{cbluebg}\else\def\fc{corangebg}\fi
      \draw[fill=\fc] (\k*\w,-0.2) rectangle ++(\w,0.4);
      \draw[fill=\fc] (\k*\w,-0.65) rectangle ++(\w,0.4);
      \node[font=\tiny\ttfamily] at (\k*\w+\w/2,0) {\char\ca};
      \node[font=\tiny\ttfamily] at (\k*\w+\w/2,-0.45) {\char\cb};
      \node[font=\tiny,text=cgray] at (\k*\w+\w/2,0.45) {\k};
    }
    \node[note] at (15*\w,-1.0) {빠른 드럼의 글자 $l = (30Q + k) \bmod 26$};
  \end{scope}
\end{tikzpicture}
\caption{주기 축 한 바퀴. 파란 26칸에서 판정하고 주황 네 칸 동안 가운데 드럼을 넘긴다. 아래 띠는 처음 두 회전에서
각 칸의 빠른 드럼 글자다. 빈 칸에서도 빠른 드럼은 계속 돌므로 둘째 회전은 E에서 시작한다.}
\label{fig:gm-cycle}
\end{figure}

맞물림 구간 316.2°--343.8°의 앞뒤에도 여유가 있다. 마지막 판정 칸($k = 25$)의 \code{JUDGE}는 지연이 최악이어도
305.5°에 끝나므로 맞물림까지 약 10.7°가 남는다. 이 마지막 칸에서 정지가 나면 모터가 가운데 드럼을 넘기기 전에 멈출
시간이다. 맞물림 뒤에서 다음 주기의 첫 접점이 패드에 닿는 355.9°까지는 12°가 남아 잠금이 자리를 잡는다
(\code{layout.md} 3.8절, 계산).

\section{위치 원판과 지시 드럼}\label{sec:gm-posdisc}

판정은 테스트 레지스터의 릴레이가 하지만, 기계가 지금 어느 위치에 있는지는 따로 알아야 한다. 정지한 위치를 화면에
보이고, 기준 모델과 대조하고, 마지막 위치에서 서고, 처음 자리로 돌아가기 위해서다.

\Ours F, M, S 세 축에 지름 56~mm의 \term{위치 원판}{position disc}이 하나씩 있다. 원판마다 슬롯 26개와 원점 슬롯
하나가 있고, 광센서 둘이 읽는다. T 센서는 그 축이 한 칸 갈 때마다 펄스를 하나 내고, I 센서는 그 축의 값이 0인 구간
동안 높다. 제어 보드가 T 펄스를 세고 I에서 0으로 맞추어 $(l, m, r)$를 안다(\ref{sec:gk-home}절). F의 원판은 빠른
드럼과 함께 계속 돌므로 T 펄스는 슬롯이 없는 빈 칸에서도 온다. 앞 절의 $l = (4m+k) \bmod 26$이 그대로 센서에 나타나는
것이다.

여기서 원점을 정하는 일이 생각보다 까다롭다. 세 원점 센서가 모두 높은 곳($l = m = r = 0$)을 처음 자리로 삼으면 될 것
같지만, 주기 축의 0번째 회전 안에서 빠른 드럼이 글자 0(A)을 보이는 칸이 \emph{둘} 있다. $k = 0$과 빈 칸 $k = 26$이다
(그림~\ref{fig:gm-cycle}의 위 띠에서 A가 두 번 나온다). 30칸이 26칸보다 많기 때문이다. 그래서 원점은 ``세 원점 센서가
높은 동안 \code{STROBE} 슬롯이 오르는 곳''으로 정한다. 슬롯은 $k = 0$에만 있으므로 하나로 정해진다.

구동부 앞판에는 \term{지시 드럼}{indicator drum} 세 개가 있다. 라인 축에 마이터 기어로 이어져 드럼과 함께 돌지만
배선은 없고, 글자가 역순으로 인쇄되어 있어 사람이 정지값을 바로 읽는다. 실물의 지시 드럼도 가운데 뱅크 오른쪽에
있었고 글자가 역순이었다~\cite{us6812report1944,sale-ringstellung}. 그 옆의 창으로 기어열이 보인다.

\paragraph{슬롯 원판.} 주기 축의 바깥 끝에는 지름 160~mm의 슬롯 원판이 있다. 반지름 72, 62, 52~mm의 트랙 셋에
\code{STROBE}, \code{JUDGE}, \code{CLEAR}의 슬롯이 26개씩 있고, 슬롯이 광센서를 지나는 동안이 곧 그 신호다. 시각이
각도로 새겨져 있으므로 기계가 느려지면 모든 구간이 시간으로 길어진다. 슬롯의 각도를 어떻게 정했는지는
\ref{sec:t-gear}절(그림~\ref{fig:t-slots})에 있다. \code{JUDGE}는 2.52°, \code{CLEAR}는 4.08° 폭이고, 가장자리는 모터
한 스텝(주기 축 0.03°)의 정수배에 있다. \code{STROBE}만은 물린 테스트 레지스터의 릴레이 지연을 재어 맞추므로,
\code{STROBE} 센서는 축 둘레로 $\pm3°$ 돌릴 수 있는 브래킷에 따로 달고, 그 트랙은 원판 두 장을 겹쳐 슬롯 폭을
0--1.5° 넓힐 수 있게 한다(추정). 슬롯 원판의 각도는 CAD의 \code{params.py} 한 곳에 있다(\ref{sec:gc-slotsteps}절).

\section{드럼 교체: 인덱스, 노브, 브레이크}\label{sec:gm-brake}

전시에서 관람객은 메뉴 카드를 보고 드럼을 직접 고르고 끼운다. 드럼의 색이 로터의 종류이고, 카드에 적힌 글자를 기준
표시에 맞추어 끼우는 것이 크립 오프셋을 거는 동작이다. 한 번의 실습에서 드럼 18개를 1.5분 안에, 곧 하나에 5초 안에
끼워야 한다. 이 요구가 드럼과 축이 만나는 부분의 모양을 정했다.

\paragraph{인덱스: 아무 글자로나, 스스로 제자리에.} 축에는 키(한 방향만 허락하는 홈)가 없다. 대신 축에 고정된
\term{인덱스 플랜지}{index flange}의 앞면에 90° V자 홈 26개가 방사 방향으로 파여 있고, 드럼 쪽 원판에는 지름 3~mm의
맞춤 핀 세 개가 있다(그림~\ref{fig:gm-index}). 홈이 26개이므로 드럼은 어느 글자 방향으로든 끼울 수 있다. 키를 두면
한 방향만 남아 오프셋을 걸 수 없다.

\Why 둥근 핀은 뾰족한 봉우리 위에 설 수 없다. 핀이 봉우리에서 조금이라도 비켜 있으면 V자의 경사면을 타고 미끄러져
가까운 홈의 바닥으로 들어간다. 그래서 드럼을 반 칸($\pm6.9°$) 안으로 대기만 하면 가장 가까운 칸에 저절로 앉는다.
홈에 앉은 핀은 두 경사면에 선으로 닿으므로 핀 세 개가 여섯 점에서 드럼의 각도, 중심, 기울기, 깊이를 모두 정한다.
처음 안은 양쪽 모두 높이 1.2~mm의 V자 톱니였는데, 톱니 끝이 평평하게 잘려 대는 각도의 59\%에서 평평한 끝끼리 닿아
걸렸다(결정 0002의 2절). 끝을 뾰족하게 하면 톱니가 너무 높아진다. 둥근 핀과 날카로운 홈의 짝은 걸리는 구간이 없다.

\begin{figure}[ht]
\centering
\begin{tikzpicture}[x=1cm,y=1cm]
  % (가) flange top view
  \begin{scope}
    \def\s{0.09}
    \fill[cgray!25,draw=cgray] (0,0) circle (17*\s);
    \fill[white,draw=cgray] (0,0) circle (4*\s);
    \foreach \i in {0,...,25} {
      \pgfmathsetmacro{\a}{90-\i*360/26}
      \draw[cgray!80!black] (\a:{10.5*\s}) -- (\a:{17*\s});
    }
    \foreach \i in {0,9,17} {
      \pgfmathsetmacro{\a}{90-\i*360/26}
      \draw[line cap=round,line width=2.6mm,cblue] (\a:{11.2*\s}) -- (\a:{16.6*\s});
    }
    \node[font=\footnotesize\bfseries] at (0,-2.3) {(가) 플랜지와 핀 셋};
    \node[note,anchor=west] at (1.75,1.0) {V홈 26개\\(13.846° 간격)};
    \node[note,anchor=west,text=cblue] at (1.75,-0.5) {핀 셋: 0번·9번·17번 칸\\(0° / 124.6° / 235.4°)};
  \end{scope}
  % (나) section: pin in a V groove
  \begin{scope}[xshift=7.9cm,yshift=-0.2cm]
    \def\u{0.55}
    \fill[cgray!25] (-5.07*\u,-2.4*\u) -- (-5.07*\u,0) -- (-3.38*\u,-1.69*\u) -- (-1.69*\u,0) -- (0,-1.69*\u) -- (1.69*\u,0) -- (3.38*\u,-1.69*\u) -- (5.07*\u,0) -- (5.07*\u,-2.4*\u) -- cycle;
    \draw[cgray!80!black] (-5.07*\u,0) -- (-3.38*\u,-1.69*\u) -- (-1.69*\u,0) -- (0,-1.69*\u) -- (1.69*\u,0) -- (3.38*\u,-1.69*\u) -- (5.07*\u,0);
    \fill[cblue!70] (0,{(-1.69+2.12)*\u}) circle (1.5*\u);
    \draw[cblue!50,dashed] (1.69*\u,{(0+1.5)*\u}) circle (1.5*\u);
    \draw[->,cred,line width=1pt] (2.45*\u,2.0*\u) to[bend left=30] (3.3*\u,0.9*\u);
    \node[note,anchor=west,text=cred,align=left] at (2.9*\u,2.6*\u) {봉우리 위에 설 수 없어\\옆 홈으로 미끄러진다};
    \node[note,anchor=north,align=center] at (0,-2.55*\u) {홈에 앉은 핀: 두 경사면에 선으로 닿는다};
    \node[font=\footnotesize\bfseries] at (0,-4.0*\u) {(나) 반지름 14~mm에서 자른 단면};
  \end{scope}
\end{tikzpicture}
\caption{인덱스. (가) 축 쪽 플랜지(지름 34~mm)의 V홈과 드럼 쪽 핀 셋의 자리. 26은 3으로 나뉘지 않으므로 핀의 간격은
124.6° / 110.8° / 124.6°다. (나) 핀(지름 3)과 90° 홈(피치 3.38, 깊이 1.69). 홈의 깊이를 반지름에 비례하게
($\pi r / 26$) 해서 봉우리가 홈의 띠 전체에서 뾰족하다.}
\label{fig:gm-index}
\end{figure}

\paragraph{노브: 4분의 1바퀴.} 드럼을 축에 붙드는 것은 드럼 앞의 노브다. 처음에는 나사(M6)였으나 4.5바퀴를 돌려야 해
하나에 약 9초가 걸렸다. 지금은 \term{베이오넷}{bayonet} 방식이다. 카메라 렌즈를 끼우듯 축 끝의 가로 핀이 노브 안의
나선 홈을 따라 90° 돌며 1.8~mm 당겨지고, 끝의 오목한 자리에서 잠긴다. 노브 안의 접시 스프링이 드럼을 약 180~N으로
누른다. 이 힘이 접점 104개가 밀어내는 힘(최대 87~N)과 인덱스가 벌어지려는 힘을 함께 이기도록 정한 값이다. 손 토크는
0.38--0.44~N·m(계산), 하나를 끼우는 데 약 4.5초(추정)다.

순서가 영리하다. 드럼을 밀면 먼저 핀이 홈에 들어가고(자리까지 1.4~mm 남았을 때), 그다음에 접점의 공이 패드에 닿는다
(0.76~mm 남았을 때). \emph{각도가 정해진 뒤에} 접점의 마찰이 생기는 것이다. 손이 미는 힘은 접점이 닿기 전까지 0이고,
그 뒤는 노브의 캠이 맡는다. 억지로 드럼을 비틀면 핀이 홈의 경사면을 타고 올라와 드럼이 약 1.5~N·m에서 한 칸 넘어간다.
부서지기 전에 넘어가고, 넘어간 것이 눈에 보인다.

\paragraph{브레이크: 덮개가 열리면 F를 잡는다.} 드럼을 바꾸는 동안 기계의 상태가 그대로여야 한다. 덮개를 열면 인터록이
모터의 48~V를 끊어 모터가 힘을 잃는다. 가운데와 느린 행은 잠금 원호가 잡고 있지만 빠른 행 F는 자유롭다. 관람객(특히
아이)이 드럼을 돌려 보면 F가 돌고, 26:30 기어로 C가 돌고, 약 1.15바퀴마다 자리올림이 가운데 행을 한 칸 넘긴다. 그 뒤에
``기준 표시에 맞춰'' 끼운 드럼은 앞서 끼운 드럼과 다른 기계 상태에서 맞춰진 것이 된다. 오프셋이 \emph{소리 없이} 틀린다
(\code{hardware/common/rotor-cartridge.md} 4.2절).

\General 엘리베이터의 브레이크는 스프링이 늘 잡고 있고, 전자석에 전기가 들어올 때만 풀린다. 정전이 되면 저절로 잡힌다.
이런 \term{무여자 작동 브레이크}{spring-applied, power-off brake}는 ``고장 나면 안전한 쪽''으로 설계하는 전형이다.

\Ours F 축에 24~V 무여자 작동 브레이크를 단다(구동부의 뱅크 쪽 칸, 바깥지름 64~mm 이하, 품번은 미확인). 브레이크
코일의 전원은 앞 덮개 스위치 두 개의 둘째 b접점을 직렬로 지난다. 덮개가 닫혀 있으면 코일에 전기가 들어가 브레이크가
풀리고, 덮개가 열리면 잡는다. F가 잡히면 26:30 기어로 C가 잡히고, M과 S는 잠금 원호가 잡는다. 기계 전체가 선다.
정지 토크는 4~N·m 이상으로 정했다. 드럼 두 개를 함께 비틀어 둘의 인덱스가 넘어가는 토크(2 × 1.5~N·m)보다 크게 해서,
억지로 비틀면 기계가 아니라 드럼의 인덱스가 넘어가게 한 것이다. 손 크랭크의 셔터 스위치는 이 고리에 넣지 않았다.
덮개를 닫은 채 크랭크로 돌려 보이는 시연은 그대로 된다. 운전 중에 덮개를 열면 브레이크가 F를 약 9~ms 안에 세운다
(계산). 48~V를 끊는 안전 기능 자체는 바뀌지 않고, 브레이크가 고장 나 풀린 채여도 브레이크가 없던 설계와 같아질
뿐이다.

관람객의 순서는 이렇게 된다(\code{layout.md} 5.2절).
\begin{enumerate}
\item \code{HOME}을 누르고 원위치 표시(\code{AT\_HOME})를 기다린다. 드럼 위의 기준 표시는 \code{HOME} 상태를 기준으로 한다.
\item 앞 덮개를 연다. 48~V가 끊기고 브레이크가 F를 잡는다.
\item 드럼마다: 노브를 4분의 1바퀴 풀어 뺀다, 카드의 색 드럼을 집는다, 카드의 글자를 기준 표시에 맞춰 민다, 노브를
4분의 1바퀴 돌린다.
\item 덮개를 닫는다. 브레이크가 풀리고 48~V가 돌아온다. 시작 스위치를 껐다 켜야 기계가 움직인다
(\ref{sec:gk-safety}절의 48~V 규칙).
\end{enumerate}

드럼을 끼우지 않았거나 노브를 돌리지 않으면 그 유닛의 길이 끊긴다. 하드웨어로 따로 검출하지 않는다. 첫 위치에서
레지스터가 정지를 내고, 감시 장치가 기준 모델의 예상과 견주어 ``드럼 장착을 확인하세요''를 띄운다.

\section{정리}

\begin{itemize}
\item 드럼은 부품 없는 4층 기판이다. 패드 104개가 네 링에 있고, 로터 배선 두 벌(가는 길 링 4$\to$3, 오는 길 링 2$\to$1)이
안쪽 층에 있다. 드럼이 글자 $o$를 보이면 접점 $j$는 패드 $j+o$에 닿아 코어의 로터 함수 $j \mapsto W[j+o]-o$가 된다.
\item 스프링 접점은 서 있는 쪽(기계)에, 패드는 도는 쪽(드럼)에 둔다. 접점은 옆으로 지나가는 용도의 구르는 공 끝이다.
\item 판정 창은 완전 접촉 구간의 가운데 절반이고, 빠른 드럼의 허용 오차는 $\pm1.1°$다. 백래시가 쌓이지 않도록 라인
축과 마이터 기어로 모든 드럼을 같은 축에서 받는다.
\item F:C = 26:30. 주기 축 한 바퀴의 30칸 가운데 26칸에서 판정하고 4칸 동안 자리올림을 한다. 빠른 드럼의 시작 글자가
주기마다 네 칸씩 밀린다.
\item 자리올림은 한 이빨 기어가 넘기고 잠금 원호가 붙든다. 둘째 자리올림은 M의 한 걸음 안에 끝나야 해서 54잇수 상당으로
크다. 디텐트는 마찰을 이기지 못하므로 달지 않는다.
\item 원점은 세 원점 센서만으로 정해지지 않는다(빈 칸 $k = 26$). \code{STROBE} 슬롯의 오름을 함께 본다.
\item 인덱스는 V홈 26개와 둥근 핀 셋, 노브는 4분의 1바퀴 베이오넷, 덮개가 열리면 무여자 작동 브레이크가 F를 잡아 드럼을
바꾸는 동안 기계의 상태를 지킨다.
\end{itemize}
